\section{A contrast: the Rovelli--Smolin operator}\label{sec:rs}

The vertices on which \cref{thm:closing} closes the gap all have an edge of spin
$\frac12$. The Rovelli--Smolin operator, in the form of De Pietri and
Rovelli~\cite[eqs.~(7.15), (7.42)]{DPR1996}, takes the absolute value of each triple
before summing and has no orientation signs:
\[
  Q_{RS}=\sum_{I<J<K\le N}|\hat q_{IJK}|,\qquad \hat V_{RS}=c_{RS}\sqrt{Q_{RS}},
\]
which is also the form Ashtekar and Lewandowski re-derive in their
appendix~\cite{AL1998} and Kisielowski studies~\cite{Kis2021}. On exactly the spin labels
where the Ashtekar--Lewandowski gap closes, it has a uniform gap.

\begin{theorem}{Rovelli--Smolin gap with a spin-\texorpdfstring{$\frac12$}{1/2} edge}{rs}
Let a five-valent vertex have nonzero spins, one of them $\frac12$, and
$\dim\Inv\ge2$. Then $Q_{RS}\ge3\sqrt3$ on the whole invariant space; in particular
$Q_{RS}$ has trivial kernel. The variant $\sum_T\sqrt{|\hat q_T|}$ of Giesel and
Thiemann~\cite{GT2006} is at least $3\cdot3^{1/4}$.
\end{theorem}

\begin{proof}
Label the spin-$\frac12$ edge $1$. The six triples containing it split into three
complementary pairs $\{1JK\},\{1LM\}$ with $\{J,K,L,M\}=\{2,3,4,5\}$.

The triple $\hat q_{1JK}$ commutes with $S=|J_L+J_M|$, and on each value of $S$ it is the
four-valent operator on $\Inv(\frac12,j_J,j_K,S)$, of dimension at most two because
$x=|J_J+J_K|\in\{S\pm\frac12\}$. On a two-dimensional block it is
$\bigl(\begin{smallmatrix}0&-4a\\4a&0\end{smallmatrix}\bigr)$, so $|\hat q_{1JK}|=4a$ there,
and by \eqref{eq:ak} with the factor $4k^2-1$ cancelled ($k=j_J+\frac12$),
$16a^2=((j_K+S+1)^2-k^2)(k^2-(j_K-S)^2)\ge k(2k+1)\ge3$. So
$|\hat q_{1JK}|\ge\sqrt3\,P_{JK}$, with $P_{JK}$ its range projector.

The kernel of $\hat q_{1JK}$ is spanned by the one-dimensional blocks, each a joint
eigenvector of $x=|J_J+J_K|$ and $y=|J_L+J_M|$ with $(x,y)$ equal to
$(j_J+j_K,\,j_J+j_K+\frac12)$ or $(|j_J-j_K|,\,|j_J-j_K|-\frac12)$. Joint eigenspaces of
$(x,y)$ are at most one-dimensional (the spin-$\frac12$ leg must couple $x$ and $y$ to
zero), so these are lines of one orthonormal basis. The kernel lines of $\hat q_{1LM}$
lie in the same basis with the roles of $x$ and $y$ exchanged. A common line would need
$j_J+j_K=j_L+j_M\pm\frac12$ simultaneously in both directions, which is impossible, or
one spin equal to the sum of the other four, which forces $\dim\Inv=1$. So the two kernel
projectors are orthogonal, $P_{JK}+P_{LM}\ge1$, and
\[
  Q_{RS}\ \ge\ \sum_{T\ni1}|\hat q_T|\ \ge\ \sqrt3\sum_{\rm pairs}(P_{JK}+P_{LM})\ \ge\ 3\sqrt3 .
\]
For the variant, $|\hat q_T|\ge\sqrt3P_T$ gives $\sqrt{|\hat q_T|}\ge3^{1/4}P_T$.
\end{proof}

The bound is not sharp: the minimum over all spin-$\frac12$ tuples with spins up to $20$
is $15.93$, at $(\frac12,\frac12,\frac12,\frac12,1)$ (and we reproduce it with our own code
on small tuples). Along the Ashtekar--Lewandowski family $(\frac12;j-\frac12,j,j,j)$ the
smallest eigenvalue of $Q_{RS}$ grows, from $21.4$ at $2j=2$ to $2709$ at $2j=40$, while
\cref{thm:closing} has the Ashtekar--Lewandowski value falling like $2^{-2j}$.

\paragraph{What the contrast does and does not say.}
The mechanism of \cref{thm:closing} is a cancellation between two triple products that
enter with signs; taking the absolute value of each triple first removes it. It is the
order of the absolute value and the sum that matters, not the presence of orientation
signs: the signless sum $|\sum_T\hat q_T|$ has exact zero eigenvalues on spin-$\frac12$
tuples too. The Rovelli--Smolin operator is not therefore gapped in general. It has
zero-volume states at valence five, for instance at five spins $2$ (\cref{sec:operator});
its general gap at valence five is open, with numerical support; and since the
consistency check of Giesel and Thiemann~\cite{GT2006}, which found it incompatible with
the quantisation of the fluxes, it has been disfavoured for dynamics, which is why
Brunnemann and Rideout study the Ashtekar--Lewandowski operator. The original paper
of Rovelli and Smolin~\cite{RS1995} (in its arXiv version, which was never revised;
we could not obtain the two-page erratum) states the higher-valence operator in a form that, read literally, takes the
absolute value after summing over triples; \cref{thm:rs} does not apply to that form,
and the later works cited here~\cite{DPR1996,AL1998,GT2006,Kis2021} use the per-triple
form above.
